\section{Week 12: Actual Decoder Semantic Refinement}

Week~12 closes \claimid{OBL-RANS-DECODE-REFINEMENT} as
\status{proved, exact-trace partial correctness}.  The target is the actual
generated \procname{RansDecodeTarget.M._rans_decode}; it is not a wrapper,
observer, or pure decoder theorem.

\subsection*{Claim boundary and theorem surface}

The compiled theorem is
\procname{actual_rans_decode_trace_refinement} in
\theoryname{Mode2RansDecoderTopHoare}.  Its precondition binds the actual
generated decoder state and buffer to a canonical exact trace:
\begin{itemize}
  \item \texttt{symsp = decoded0};
  \item \texttt{statep = state0};
  \item \texttt{bufp = buffer0};
  \item the literal mode-2 \texttt{symbol\_words} and \texttt{dsyms\_words}
        tables;
  \item \texttt{actual\_mode2\_decoder\_trace\_input expected\_symbols
        buffer0 state0 encoded\_size}.
\end{itemize}
The exact-trace premise contains the canonical 1024-symbol stream, the exact
\texttt{traceBytes} buffer, the encoded-size relation
\(\mathit{encoded\_size} = |\mathsf{traceBytes}(\mathsf{symbolListOfArray}
(expected\_symbols))|\), and the concrete state fields
\((count,size,m)=(1024,encoded\_size,13)\).  It does \emph{not} contain
decoder success, output equality, or a post-state observation.

The postcondition published by the theorem is:
\[
\begin{aligned}
  bad=0
  &\;\land\; off=\mathit{encoded\_size} \\
  &\;\land\; decoded[0..1024)=expected\_symbols[0..1024) \\
  &\;\land\; decoded[1024..2048)=decoded_0[1024..2048).
\end{aligned}
\]
The internal final-state fact \(x=2^{23}\) is proved at the generated loop
exit and used to discharge the actual final reject checks, but it is not
exported as an ABI result.

\subsection*{Reference state, cursor, and 0/1/2-byte normalization}

For \(syms=\mathsf{symbolListOfArray}(expected\_symbols)\), the reference
decoder boundary is given by:
\[
  state(i)=\mathsf{decoderStateAt}(syms,i),\qquad
  cursor(i)=\mathsf{decoderCursor}(syms,i).
\]
The compiled cursor lemmas establish
\[
\begin{aligned}
  cursor(0)&=4,\\
  cursor(i+1)&=cursor(i)+|\mathsf{segment}(i)|,\\
  cursor(1024)&=|\mathsf{traceBytes}(syms)|,
\end{aligned}
\]
and \(state(0)\) is the parsed little-endian head while \(state(1024)=2^{23}\).
The concrete mode-2 bound keeps each normalization segment at length
0, 1, or 2 bytes.

The actual inner invariant tracks the partial replay cursor
\(k=off-cursor(i)\) and the current reduced state.  The generated guard
matches the pure replay relation: when one more segment byte remains, the loop
consumes it; otherwise the no-consume exit establishes the next state/cursor
boundary.  This is the exact 0/1/2-byte case split used by the theorem
\procname{decoder\_inner\_trace\_step}, \procname{decoder\_inner\_trace\_stop},
and \procname{decoder\_inner\_trace\_exit\_to\_outer}.

\subsection*{Concrete table bridge}

The word-step bridge is split from the loop proof.  The generated code uses the
packed \texttt{symbol\_words} lookup on \(x \mathbin{\&} 1023\), stores the
recovered symbol, loads the packed \texttt{dsyms\_words} entry, and computes
the reduced state with the exact mode-2 arithmetic.  The compiled lemma family
\breakablett{generated_decoder_lookup_word_from_mode2},
\breakablett{generated_decoder_symbol_from_mode2}, and
\breakablett{generated_decoder_word_update_from_mode2} connects those exact
generated expressions to the mathematical decoder step without introducing a
table axiom.

\subsection*{Exact exit and final checks}

At outer-loop exit, \procname{decoder\_outer\_trace\_exit\_components} proves
\[
  i=1024,\qquad x=2^{23},\qquad off=\mathit{encoded\_size},\qquad
  bad=0,
\]
together with the full decoded-prefix equality and the untouched
\([1024,2048)\) output frame.  Those facts discharge the generated final-state
parse and reject checks.  The theorem is therefore exact-trace partial
correctness, not losslessness or termination.

\subsection*{Non-vacuity and remaining edge}

The exact-trace premise is satisfiable as a pure trace relation and is the
mathematical relation produced by the Week~11 encoder success theorem.  It
does not prove that the actual encoder success branch is reachable.
\claimid{OBL-RANS-ACTUAL-SUCCESS-WITNESS} therefore remains
\status{partial}.

The remaining implementation edge is now only the unary harness composition:
transport the encoder suffix and the actual copy-helper slice facts into the
decoder's exact pointwise trace-read premise.  Until that theorem compiles,
\claimid{OBL-RANS-CORE-INVERSE} and the HBZ parent remain
\status{partial}.

\subsection*{Rejected tactics and follow-up}

A monolithic \texttt{inline; sim} proof was rejected because the outer symbol
loop and the variable-length normalization loop require different invariants.
Extending the older control theorem was also rejected: it can only retain
\(\texttt{bad}\in\{0,1\}\) and cannot express symbol recovery.  The Week~13
single edge is the actual encoder\(\rightarrow\)copy\(\rightarrow\)decoder
harness theorem; no new codec layer should be introduced before it compiles.
